% !TEX TS-program = pdflatex
% CV réseau LinkedIn — Écosystème SAP (partenaires, ex-collègues Streamlink)
\documentclass[a4paper,9pt]{article}
\usepackage[utf8]{inputenc}
\usepackage[T1]{fontenc}
\usepackage[scaled=0.90]{helvet}
\renewcommand{\familydefault}{\sfdefault}
\usepackage[left=0.8cm,right=0.8cm,top=0.3cm,bottom=0.25cm]{geometry}
\usepackage{enumitem}
\usepackage{titlesec}
\usepackage{xcolor}
\usepackage[hidelinks]{hyperref}
\definecolor{navy}{RGB}{23,54,93}
\pagestyle{empty}
\setlength{\parindent}{0pt}
\setlength{\parskip}{0pt}
\linespread{0.82}
\hypersetup{pdftitle={CV - Baha Eddine Belhaj Mouldi - SAP Basis Security},pdfauthor={Baha Eddine Belhaj Mouldi},pdfsubject={SAP Basis, Securite SAP, Detection de menaces, S4HANA}}
\titleformat{\section}{\normalsize\bfseries\color{navy}}{}{0pt}{\MakeUppercase}[{\vspace{1pt}\color{navy}\hrule height 0.6pt}]
\titlespacing*{\section}{0pt}{1.5pt}{0.8pt}
\setlist[itemize]{leftmargin=1.0em,label=\textbullet,itemsep=0pt,topsep=0pt,parsep=0pt}
\newcommand{\cventry}[4]{\textbf{#1} \hfill \textbf{#4}\\[0.2pt]#2{\ifx\relax#3\relax\else, #3\fi}\par\vspace{0.3pt}}
\begin{document}
\begin{center}
{\LARGE\bfseries\color{navy} Baha Eddine Belhaj Mouldi}\\[2pt]
{\normalsize\color{navy} SAP BASIS \& Sécurité \textbar\ Détection de Menaces \textbar\ S/4HANA \textbar\ Automatisation}\\[3pt]
{\small Tunis, Tunisie \quad|\quad +216 55 128 982 \quad|\quad \href{mailto:bahaeddine.belhajmouldi@gmail.com}{bahaeddine.belhajmouldi@gmail.com}}\\[1pt]
{\small \href{https://www.linkedin.com/in/bahamouldi}{linkedin.com/in/bahamouldi} \quad|\quad \href{https://github.com/bahamouldi}{github.com/bahamouldi} \quad|\quad \href{https://bahamouldi.github.io/portfolio/}{bahamouldi.github.io/portfolio}}
\end{center}
\vspace{1pt}
\section{Profil}
Ingénieur en génie informatique (ENIT, mention Très Bien) avec une expérience concrète de détection de menaces sur environnement SAP, acquise lors d'un stage chez Streamlink. Pratique de l'administration SAP BASIS, de la surveillance de transactions critiques (SU01, SE16N, PFCG) et de la réponse automatisée à incident (blocage via BAPI\_USER\_LOCK). Complément par une solide expérience de sécurité applicative (BeeWAF, conformité 7 référentiels) et d'automatisation Python. Rigueur, sens de la confidentialité et goût pour les environnements ERP exigeants.
\section{Compétences clés}
\textbf{SAP} : Administration SAP BASIS, environnement S/4HANA, transactions critiques (SU01, SE16N, PFCG), séparation des tâches (SoD), PyRFC\\[0.3pt]
\textbf{Détection \& réponse à incident} : pipeline automatisé (Python, ELK, N8N, TheHive), scoring de risque, blocage automatisé (BAPI\_USER\_LOCK)\\[0.3pt]
\textbf{Sécurité applicative} : OWASP Top 10, gestion des correctifs (CVE), conformité (ISO 27001, GDPR, PCI DSS, SOC 2)\\[0.3pt]
\textbf{Systèmes \& Réseaux} : Windows Server, Active Directory (GPO), Linux, VPN, Proxy, Firewall\\[0.3pt]
\textbf{Automatisation} : Python, Bash, PowerShell, Docker, Kubernetes\\[0.3pt]
\textbf{Supervision} : Elastic Stack, Kibana, Metricbeat, alerting temps réel
\section{Expérience professionnelle}
\cventry{Stagiaire Cybersécurité --- Détection \& Réponse à incident SAP}{Streamlink}{Tunisie}{07/2025 -- 08/2025}
\begin{itemize}
  \item Pipeline de détection automatisé : collecte Python/PyRFC, corrélation ELK, orchestration N8N, suivi TheHive. Détection $<$2 min, blocage automatisé $<$7 min.
  \item Surveillance de 7 transactions critiques SAP (SU01, SE16N, PFCG) avec scoring de risque et alerting temps réel via Kibana.
\end{itemize}
\cventry{Ingénieur Sécurité --- Stagiaire PFE (mention Très Bien)}{Beehive Entreprises}{Tunisie}{02/2026 -- 06/2026}
\begin{itemize}
  \item Conception de BeeWAF : conformité continue sur 7 référentiels, gestion des correctifs (37 CVE), déploiement Kubernetes en production.
\end{itemize}
\cventry{Spécialiste Cybersécurité, Freelance (1 an)}{Beehive Entreprises}{Tunisie}{01/2025 -- 12/2025}
\begin{itemize}
  \item Audits de vulnérabilités sur 8 applications web et API, rapports de remédiation.
\end{itemize}
\cventry{Chef de projet technique, Indépendant --- Équipe de 8 personnes}{VMate}{Tunisie}{01/2024 -- 06/2024}
\begin{itemize}
  \item Coordination d'équipe, reporting.
\end{itemize}
\section{Projets clés}
\cventry{Système de Détection et Réponse aux Incidents SAP}{Python, PyRFC, ELK, N8N, TheHive}{}{2025}
\begin{itemize}
  \item Pipeline complet : logs SAP $\to$ Metricbeat $\to$ Elasticsearch $\to$ Kibana $\to$ N8N $\to$ TheHive. Réponse automatisée sur 3 scénarios d'incident (connexions anormales, escalade de privilèges, accès non autorisé).
\end{itemize}
\cventry{BeeWAF --- Sécurité Applicative en Production (PFE, 86/100)}{FastAPI, Docker, Kubernetes, ELK}{}{2026}
\begin{itemize}
  \item Conformité continue sur 7 référentiels, gestion structurée des correctifs.
\end{itemize}
\section{Formation}
\cventry{Diplôme d'Ingénieur en Génie Informatique --- Mention Très Bien}{École Nationale d'Ingénieurs de Tunis (ENIT)}{Tunisie}{09/2023 -- 06/2026}
\begin{itemize}
  \item Diplômé le 25/06/2026. Classé 65e sur 600+ au concours national d'entrée.
\end{itemize}
\cventry{Classes préparatoires Physique-Technologie}{Institut Préparatoire aux Études d'Ingénieurs, El Manar}{Tunisie}{09/2021 -- 06/2023}
\section{Certifications}
Réseaux (CCNA) --- Cisco (2025) ; Fondements du Deep Learning --- NVIDIA (2025) ; Machine Learning --- NVIDIA (2024) ; Git \& GitHub --- 365 Data Science (2024).
\section{Langues}
Français (Courant) \quad|\quad Anglais (B2) \quad|\quad Arabe (Natif)
\end{document}
